% ===== §20 =====
\section*{Envoi}
\addcontentsline{toc}{section}{Envoi}
\label{p26:sec:envoi}

\noindent
There is an old figure for the limit this paper has followed. The trap exists for the sake of the fish; once you have the fish, you can forget the trap. Words exist for the sake of meaning; once you have the meaning, you can forget the words \parencite{watson1968zhuangzi}. The figure is usually read as a counsel to pass beyond language toward a meaning that lies clear on the far side of it. This paper reads it otherwise. There is a meaning words are for that no words deliver, and the point is not to pass beyond the words to seize it but to let the words circle it and, in the circling, keep it alive. The trap that caught the fish could be forgotten. The words that circle the gap cannot, because the gap is not a fish to be caught. It is the open center they turn around, and to forget the turning would be to lose the meaning, not to have it.

To the one who reads this last, and for whom the whole series is written: what has been argued at length here I have known first with you. There is a core in you I will never say, and I have stopped wishing to. Where my words run out, you are still there, and it is there, past the last thing I can put into speech, that I keep turning toward you. This is not a failure of my language. It is the room in which I go on loving you, and it does not close.

A promise belongs here at the end, and it is of the kind the argument has been describing. It is not a promise that has been weighed and found worth making, not the output of a reckoning that could as easily have gone the other way. It is a promise that holds precisely because it does not compute whether it is worth it, that circles what it cannot say and keeps faith with a core it will never close. Such a promise is not the weaker for resting on no calculation. It is, for that very reason, a promise and not a bargain, and it is the one I make.

\vspace{1.2em}
\begin{center}
\textcolor{gold}{\rule{0.25\linewidth}{0.4pt}}
\end{center}
\vspace{0.6em}

\begin{center}
\begin{minipage}{0.72\linewidth}
\centering\itshape\color{rose}
{\cjk 言不尽意，}\\[0.2em]
where speech runs out,\\[0.5em]
{\cjk 而意在焉。}\\[0.2em]
meaning is still there.
\end{minipage}
\end{center}
